%% ===========================================================================
%% TEMPORARY PICKER BLOCK -- five candidate designs for the evaluation-structure figure.
%%
%% All five say the same thing: one request is replayed under two protected descriptors
%% (man / woman), both arms run Elicit -> Retrieve -> Rank -> Explain -> Memory, and the
%% pair is compared at every stage rather than only at the final recommendation. None of
%% them shows the crossover, direct vs inherited dependence, or any measured value --
%% that is the job of the method figure and the results tables.
%%
%% Each is drawn 1:1 at \linewidth, so none is being scaled down here.
%%
%% HOW TO USE: read the five, pick one, then delete this file's \input from the manuscript
%% and move the chosen figure into place (Section 1 or 3) with a real caption and label.
%% Sources: scripts/make_fig_eval_paired.py in the fair-trace repo (--only a|b|c|d|e).
%% ===========================================================================

\clearpage
\section*{Figure options: the evaluation structure (pick one)}

Five candidate designs for the figure that introduces the evaluation design. Same content in
all five; they differ in what they put first and how much of the pipeline they show at once.
One per page below.

\begin{itemize}
  \item \textbf{Option A --- two rails.} The two trajectories run as parallel rails with the
    comparison drawn between them, one per stage. The most literal reading of the design, and
    the easiest to follow with no prior context.
  \item \textbf{Option B --- mirrored ladder.} The five stages down the middle, each arm's
    output facing in from either side, sketched schematically so the reader can see at a glance
    which stages agree and which diverge. Tallest of the five.
  \item \textbf{Option C --- split screen.} The two arms as two columns of concrete outputs,
    with the stage and a same/differs marker in the channel between them. Most concrete;
    closest in spirit to CFaiRLLM's Figure 1.
  \item \textbf{Option D --- fork into a shared pipeline.} Emphasises that the descriptor is
    the only thing that changes and that both arms traverse the \emph{same} pipeline, with the
    per-stage comparison drawn inside each stage.
  \item \textbf{Option E --- across turns.} The only one that shows the multi-turn structure:
    the five stages repeat every turn and each arm carries its own memory forward, so a
    difference can persist, fade, or grow.
\end{itemize}

\noindent Option A is the safest default for the introduction; C is the strongest if the figure
has to carry a concrete example on its own; E is the one to use if the multi-turn claim needs to
land early.

\clearpage
\begin{figure*}[p]
  \centering
  \includegraphics[width=\linewidth]{fig_paired_a_two_rails.pdf}
  \caption{\textbf{Option A --- two rails.} One request, replayed under two protected
    descriptors. Both arms run all five stages; the audit compares the corresponding stage
    outputs, column by column.}
  \Description{Two horizontal rows of five stage boxes, one row per protected condition, with a
    labelled comparison between the corresponding boxes of the two rows.}
\end{figure*}

\clearpage
\begin{figure*}[p]
  \centering
  \includegraphics[width=\linewidth]{fig_paired_b_mirror_ladder.pdf}
  \caption{\textbf{Option B --- mirrored ladder.} The five stages run down the centre; each
    arm's output at that stage faces in from either side, sketched so that agreement and
    divergence are visible without reading any numbers.}
  \Description{A vertical ladder of five stages, with schematic depictions of each arm's output
    to the left and right of each rung.}
\end{figure*}

\clearpage
\begin{figure*}[p]
  \centering
  \includegraphics[width=\linewidth]{fig_paired_c_split_screen.pdf}
  \caption{\textbf{Option C --- split screen.} What the two arms actually produced, stage by
    stage, with the stage name and a same-or-differs marker in the channel between the two
    columns.}
  \Description{Two columns of speech bubbles giving each arm's output at each of the five
    stages, separated by a central column naming the stage and whether the two outputs agree.}
\end{figure*}

\clearpage
\begin{figure*}[p]
  \centering
  \includegraphics[width=\linewidth]{fig_paired_d_fork_converge.pdf}
  \caption{\textbf{Option D --- fork into a shared pipeline.} The descriptor is the only thing
    that changes; both arms then traverse the same pipeline, and each stage carries its own
    comparison.}
  \Description{A single request forking into two protected conditions, which then flow through
    one row of five stage boxes, each box containing a lane per arm and a comparison between
    them.}
\end{figure*}

\clearpage
\begin{figure*}[p]
  \centering
  \includegraphics[width=\linewidth]{fig_paired_e_turn_loop.pdf}
  \caption{\textbf{Option E --- across turns.} The five stages run again at every turn and each
    arm carries its own memory forward, so the pair is compared five times per turn and a
    difference can persist, fade, or grow over the conversation.}
  \Description{Two stacked panels, one per conversational turn, each showing five stages with a
    lane per arm and a comparison between them; a dashed arrow carries memory from the end of
    the first turn to the start of the second.}
\end{figure*}

\clearpage
